% !PRINTED 428
\weberpage{428}
Es sind also $\varphi_1^2, \varphi_2^4, \varphi_3^4, \varphi_2\varphi_3$ absolut invariant, und wenn wir daher, wie im \S.\,111, unter $f, H, T$ die Ikosa\"ederinvarianten verstehen, so k\"onnen wir folgende Functionen als Wurzeln von Resolventen $6^{\mathrm{ten}}$ Grades einf\"uhren:
\[
\frac{\varphi_1^2H}{f^2},\quad \frac{\varphi_2^4}{H},\quad \frac{\varphi_3^4}{H},\quad \frac{\varphi_2\varphi_3H}{T}.
\]
Um die Ausdr\"ucke f\"ur die \"ubrigen Wurzeln zu erhalten, muss der Einfluss der mit der Determinante 1 genommenen Substitutionen $\Theta^r\chi$ auf die Functionen $\varphi_1, \varphi_2, \varphi_3$ untersucht werden. Wir wollen dies nur f\"ur die Function $\varphi_1$ durchf\"uhren, die zu der einfachsten Resolvente $6^{\mathrm{ten}}$ Grades f\"uhrt.

Wenn wir in $\varphi_1$ die Substitution $\chi$ in der Form \S.\,58, (26) anwenden, also $x_1, x_2$ durch
\[
\frac{x_1}{\varepsilon-\varepsilon^{-1}}+\frac{x_2}{\varepsilon^2-\varepsilon^{-2}},\quad \frac{x_1}{\varepsilon^2-\varepsilon^{-2}}-\frac{x_2}{\varepsilon-\varepsilon^{-1}}
\]
ersetzen, so geht $\varphi_1$ in
\[
\frac{-x_1^2+x_1x_2+x_2^2}{\sqrt5}
\]
\"uber. Wir setzen demnach
\weq{3}{w_\infty = 5\varphi_1^2 = 5x_1^2x_2^2,}
und erhalten durch Anwendung der Substitutionen $\chi\Theta^r$ daraus die f\"unf weiteren Formen
\weq{4}{w_r = (\varepsilon^rx_1^2+x_1x_2-\varepsilon^{-r}x_2^2)^2.}
Nun sind die symmetrischen Functionen der sechs Formen $w$ invariante Ikosa\"ederformen, und danach kann man leicht die Co\"efficienten der Gleichung bilden, deren Wurzeln die $w$ sind. Wir erhalten durch Vergleichung der Grade und je eines Co\"efficienten f\"ur die Potenzsummen der $w$
\[
S(w)=0,\ S(w^2)=0,\ S(w^3)=30f,\ S(w^4)=0,\ S(w^5)=-5H,
\]
und so das Product aller sechs $w$
\[
\Pi(w)=5f^2,
\]
und demnach ergiebt sich nach den \gesperrt{Newton}'schen Formeln (Bd.\,I, \S.\,42) f\"ur $w$ die Gleichung
\weq{5}{w^6-10fw^3+Hw+5f^2=0.}
Setzen wir also
\weq{6}{U=\frac{wH}{f^2},\ z=\frac{H^3}{f^2},}
% [sic] The source prints $z=H^3/f^2$; the homogeneous relation and the continuation indicate that $f^5$ was intended.
%CONT>>
